\item \model{MiniMax-M2}

\paragraph*{رد توجه پنجره لغزان به نفع توجه کامل مجهز به \en{RoPE}}
یکی از مهم‌ترین دستاوردهای تجربی گزارش‌شده در توسعه مدل \model{MiniMax-M2} \cite{minimax2026m2} با زمینه بومی ۱۹۲K توکن، ارزیابی نقادانه مکانیزم‌های توجه کارآمد در برابر توجه کامل است. بر خلاف مدل پیشین این شرکت (\en{MiniMax-Text-01}) که از توجه خطی لایتنینگ (\en{Lightning Attention}) استفاده می‌کرد، و برخلاف الگوهای پنجره لغزان (\en{SWA})، مدل \model{MiniMax-M2} از \textbf{توجه کامل (\en{Full Attention}) مجهز به \en{RoPE} در تمامی ۶۲ لایه} استفاده می‌کند.

\paragraph*{تحلیل نقایص تئوریک پنجره‌های برشی در استدلال پیچیده}
آزمایش‌های گسترده ابلیشن نشان دادند که رویکردهای هیبریدی و پنجره لغزان، اگرچه در بافت‌های کمتر از ۳۲K توکن عملکردی نزدیک به توجه کامل دارند، اما در بافت‌های فراتر از ۳۲K در آزمون‌های استدلال چندمرحله‌ای (\en{Multi-hop Reasoning}) و بازیابی پیوسته دچار سقوط معنادار می‌شوند:
\begin{itemize}
    \item در آزمون \en{RULER 128K CWE}، مدل مجهز به SWA به امتیاز ۷۲.۰ تنزل یافته در حالی که مدل مجهز به توجه کامل با RoPE امتیاز ۹۰.۰ را ثبت کرده است (جدول ۲ مقاله).
    \item در آزمون \en{SWE-verified} پس از فاز تنظیم ناظر (\en{SFT})، مدل توجه کامل نمره ۵۴.۷ را در مقایسه با ۵۰.۲ مدل SWA به دست آورده است.
\end{itemize}
بنابراین، مدل بر مبنای اعمال پیوسته ماتریس دوران $R_{\Theta, m}^d$ در تمام بردارهای پرس‌وجو و کلید در تمام طول پنجره ۱۹۲K آموزش دیده است. فرآیند توسعه طول بافت در سه مرحله منسجم از ۸K به ۳۲K و در نهایت ۱۹۲K توکن هدایت شد تا توزیع فازهای دورانی به تدریج همگرا گردد.